\subsection*{5.2 Bounded-Relativization Barriers for Explicit Constructions}
\subsubsection*{5.2.1 Barriers for Almost-Everywhere Pseudodeterministic Constructions}
We show that $\mathsf{EXPH}$-relativizing techniques cannot prove almost-everywhere pseudodeterministic constructions, even if the construction algorithms are allowed $2^{o(n)}$ time. The underlying oracle uses the query complexity lower bounds proved by Goldwasser, Impagliazzo, Pitassi, and Santhanam [31] in a black-box fashion. (We thank Rahul Santhanam for pointing out that the query complexity lower bounds imply an $\mathsf{EXPH}$-computable oracle without almost-everywhere pseudodeterministic constructions in a black-box fashion.)
\paragraph{Definition 50.} A pseudodeterministic decision tree for a search problem $S$ is a distribution $\mathcal T$ over decision trees with the following property: For every input $x$, there is a canonical value $o$ such that with probability at least $2/3$, $\mathcal T(x)=o$. Let $\operatorname{psPdt}(S)$ denote the minimum depth of any pseudodeterministic decision tree for $S$.
Consider the following search problem denoted as $\mathsf{Find1}$. The input is a string $x\in\{0,1\}^N$ where it is guaranteed that there are at least $N/2$ bits in $x$ that are equal to $1$. The problem is to find some $i\in[N]$ such that $x_i=1$. The following lower bound on the pseudodeterministic query complexity of $\mathsf{Find1}$ was proved in [31]:
\paragraph{Theorem 51 [31].} $\operatorname{psPdt}(\mathsf{Find1})=\Omega(\sqrt N)$.
We use this result as a black box and construct an oracle $\mathcal O\in\mathsf{EXPH}$ without almost-everywhere pseudodeterministic constructions running in $2^{o(n)}$ time:
